\section*{Hoofdstuk \ref{ch:b3:planetary-atmospheres} --- Planetaire atmosferen en bewoonbaarheid}

\begin{solution}{exo:b3:planetary-atmospheres:1}
$H = RT/Mg$: Venus $8.314\times737/(0.04401\times8.87) = \qty{15.7}{km}$;
de Aarde $8.314\times288/(0.02897\times9.80) = \qty{8.4}{km}$; Mars
$8.314\times214/(0.04334\times3.71) = \qty{11.1}{km}$; Titan
$8.314\times94/(0.0280\times1.35) = \qty{20.7}{km}$. De atmosfeer van Titan,
hoe koud ook, is de meest uitgebreide: zijn zwaartekracht is
zeven keer zwakker dan die van de Aarde.
\end{solution}

\begin{solution}{exo:b3:planetary-atmospheres:2}
$\lambda = GMm/kTr = \num{6.674e-11}\times\num{5.97e24}\times\num{1.674e-27}
/(\num{1.381e-23}\,T\times\num{6.878e6})$: 8.78 bij \qty{800}{K}, 4.39 bij
\qty{1600}{K}. De ontsnappende fracties zijn $9.78\,e^{-8.78} =
\num{1.5e-3}$ en $5.39\,e^{-4.39} = \num{6.7e-2}$: de zonnecyclus
verandert de ontsnappende fractie van waterstof met een factor van ongeveer 44.
\end{solution}

\begin{solution}{exo:b3:planetary-atmospheres:3}
(a)~$d = S_\mathrm{eff}^{-1/2}$: \omterm{def:b3:planetary-atmospheres:runaway}{vochtig broeikaseffect} \qty{0.993}{au}, op hol geslagen
\qty{0.975}{au}, maximaal broeikaseffect \qty{1.705}{au}. (b)~Venus, op 1.91,
ontvangt meer dan de flux bij de \omterm{def:b3:planetary-atmospheres:runaway}{limiet van het op hol geslagen broeikaseffect} en zelfs meer dan de empirische
limiet voor recent Venus (1.78): het ligt binnen de binnenrand. Mars, op 0.43,
ligt tussen 1.014 en 0.3438, binnen de conservatieve bewoonbare zone:
zijn droogte komt van zijn geringe grootte en het verlies van zijn atmosfeer, niet
van zijn afstand.
\end{solution}

\begin{solution}{exo:b3:planetary-atmospheres:4}
(a)~$\log_{10}p = 3.9895 - 443.028/93.66 = -0.741$, $p =
\qty{0.18}{bar}$. (b)~$0.05\times1.6 = \qty{0.080}{bar}$, een relatieve
luchtvochtigheid van $0.080/0.18 = 44\,\%$. (c)~De lucht boven de zeeën is ver van
verzadigd: methaan verdampt uit de zeeën, net zoals water uit de
oceanen van de Aarde verdampt, en de vochtige lucht moet stijgen en afkoelen voordat het regent.
\end{solution}

\begin{solution}{exo:b3:planetary-atmospheres:5}
$\lambda_\mathrm D = 7.02\times2.014/1.008 = 14.0$. $u_\mathrm D/u_\mathrm H
= \sqrt{1.008/2.014} = 0.71$. Verhouding per atoom $= 0.71\times(15.0\,
e^{-14.0})/(8.02\,e^{-7.02}) = 0.71\times\num{1.69e-3} = \num{1.2e-3}$.
\omterm{def:b3:planetary-atmospheres:jeans-escape}{Jeans-ontsnapping} laat bijna geen deuterium ontsnappen: het is een sterk filter, met
$\alpha \approx 0.001$.
\end{solution}

\begin{solution}{exo:b3:planetary-atmospheres:6}
(a)~$0.016/\num{1.5575e-4} = 103$. (b)~Met $\alpha = 0$ is de verrijking
$M_0/M$: $M_0 = 103\times\num{6e15} = \qty{6.2e17}{kg}$, of
$\num{6.2e17}/(1000\times\num{4.60e14}) = \qty{1.3}{m}$ water boven
Venus. (c)~$M_0/M = 103^{1/0.2} = \num{1.2e10}$, een vroegere voorraad van
\qty{7e25}{kg}, twaalf keer de massa van de Aarde: onmogelijk. De
verrijking vereist een sterk fractionerende ontsnapping (Jeans of niet-thermisch),
die ten minste tijdens het verlies van het laatste water actief is, of een
stationaire toestand met een aanvoer van water.
\end{solution}

\begin{solution}{exo:b3:planetary-atmospheres:7}
(a)~In vier miljard jaar ($\num{1.26e17}$~s) verliest Mars
$\num{3.8e43}$ atomen, $\num{1.9e43}$ watermoleculen, $\qty{5.7e17}{kg}$;
over $4\pi(\num{3.3895e6})^2 = \qty{1.44e14}{m^2}$, wat neerkomt op een laag
\qty{3.9}{m} diep. (b)~Twintig keer meer, \qty{78}{m}. (c)~Zelfs de
snelheid bij grote scheefstand bereikt alleen het onderste eind van de 100 tot
\qty{1500}{m}; de rest ging de korst en het ijs in, of ging
sneller verloren toen de Zon jong was.
\end{solution}

\begin{solution}{exo:b3:planetary-atmospheres:8}
(a)~$m = pA/g = 636\times\num{1.44e14}/3.71 = \qty{2.5e16}{kg}$.
(b)~$0.1\times\num{1.26e17} = \qty{1.3e16}{kg}$, de helft van de huidige
atmosfeer; het geheel ervan in $\num{2.5e16}/0.1 = \qty{2.5e17}{s}$,
\num{7.8}~miljard~jaar. (c)~De jonge Zon was actiever en haar
stormen kwamen vaker voor, en het verlies groeit tijdens stormen; andere
kanalen (waterstof, sputtering, fotochemische ontsnapping) worden niet meegerekend.
\end{solution}

\begin{solution}{exo:b3:planetary-atmospheres:9}
(a)~$S(1 - 0.193)/4 = 291$ geeft $S = \qty{1442}{W.m^{-2}}$, 1.06 keer
die van de Aarde. (b)~Venus zou vandaag $0.807\times2601/4 =
\qty{525}{W.m^{-2}}$ absorberen; een jong Venus, $0.7\times2601 = \qty{1821}{W.m^{-2}}$
ontvangen, \qty{367}{W.m^{-2}} geabsorbeerd: beide liggen boven de limiet, dus een
wolkenvrij jong Venus zou vanaf het begin op hol zijn geslagen. (c)~Wolken
verhogen het albedo; driedimensionale modellen met wolken hebben een pad
naar een bewoonbaar jong Venus gevonden.
\end{solution}

\begin{solution}{exo:b3:planetary-atmospheres:10}
(a)~$H = kT/mg = \num{1.381e-23}\times1000/(16\times\num{1.661e-27}\times9.21)
= \qty{56}{km}$. (b)~$n_\mathrm{exo} = 1/(\sigma_c H) = 1/(\num{3e-19}\times
\num{5.64e4}) = \qty{5.9e13}{m^{-3}}$; $z = 200 + 56.4\ln(\num{1e16}/
\num{5.9e13}) = 200 + 56.4\times5.13 = \qty{490}{km}$. (c)~Doordat de zwaartekracht hoger zwakker is,
neemt de schaalhoogte met de hoogte toe, zodat de dichtheid langzamer daalt
dan aangenomen: de \omterm{def:b3:planetary-atmospheres:exobase}{exobasis} schuift omhoog.
\end{solution}

\begin{solution}{exo:b3:planetary-atmospheres:11}
(a)~$\Phi = \num{3.0e13}\times\num{3e-3} = \qty{9e10}{cm^{-2}.s^{-1}} =
\qty{9e14}{m^{-2}.s^{-1}}$; over $4\pi(\num{6.371e6})^2 =
\qty{5.10e14}{m^2}$, $\num{4.6e29}$ moleculen $\ce{H2}$ per seconde.
(b)~Een oceaan bevat $\num{1.4e24}/18.015\times\num{6.022e23} =
\num{4.7e46}$ watermoleculen, één $\ce{H2}$ elk: het duurt
$\qty{1.0e17}{s} = \num{3.2}\text{ miljard jaar}$, vergelijkbaar met de ouderdom van
de Aarde, en daarom markeert een mengverhouding van \num{3e-3} de limiet voor het \omterm{def:b3:planetary-atmospheres:runaway}{vochtig
broeikaseffect}.
\end{solution}

\begin{solution}{exo:b3:planetary-atmospheres:12}
(a)~$p_1 = g/D\kappa = 9.81/(1.66\times0.005) = \qty{1180}{Pa}$.
(b)~$1/T = 1/273.16 - \ln(1182/611.657)\times461.5/\num{2.5e6}$, $T =
\qty{282.5}{K}$, en $\sigma T^4 = \qty{361}{W.m^{-2}}$. (c)~De werkelijke limiet
is lager, dus het stralingsniveau ligt kouder en op lagere druk: echte
waterdamp is ondoorzichtiger dan dit grijze gas, en zijn effectieve $\kappa$
groter.
\end{solution}

\begin{solution}{pb:b3:planetary-atmospheres:1}
\textbf{1.} $g = GM/R^2 = \num{6.674e-11}\times\num{5.2e24}/(\num{6.2e6})^2 =
\qty{9.03}{m.s^{-2}}$; $\bar\rho = M/\tfrac43\pi R^3 = \qty{5210}{kg.m^{-3}}$.

\textbf{2.} $v_\mathrm{esc} = \sqrt{2GM/r} = \sqrt{2\times\num{3.47e14}/\num{6.7e6}}
= \qty{10.18}{km.s^{-1}}$.

\textbf{3.} $u = \sqrt{2kT/m}$: waterstof \qty{4.81}{km.s^{-1}}, helium
\qty{2.41}{km.s^{-1}}, zuurstof \qty{1.21}{km.s^{-1}}.

\textbf{4.} $\lambda = (v_\mathrm{esc}/u)^2$: 4.49, 17.8, 71.2.

\textbf{5.} Alleen waterstof: de ontsnappende fractie is 0.06, die van helium
\num{3.5e-7}, die van zuurstof \num{9e-30}. Geen van deze ligt in de buurt van het hydrodynamische regime,
$\lambda \approx 1.5$ tot 3.

\textbf{6.} $\lambda = 4.49\times1400/3000 = 2.09$: binnen de overgang;
waterstof zou hydrodynamisch zijn weggestroomd en mogelijk
zwaardere gassen met zich mee hebben gesleept.

\textbf{7.} $(1 + 4.49)\,e^{-4.49} = 0.0618$.

\textbf{8.} $n u/2\sqrt\pi = \num{2e11}\times4806/3.545 = \qty{2.71e14}{m^{-2}.s^{-1}}$;
$\Phi_\mathrm J = \num{2.71e14}\times0.0618 = \qty{1.68e13}{m^{-2}.s^{-1}}$.

\textbf{9.} $4\pi(\num{6.7e6})^2 = \qty{5.64e14}{m^2}$; $\num{9.5e27}$ atomen per
seconde.

\textbf{10.} $\num{1.4e24}/18.015 = \qty{7.8e22}{mol}$, $\num{4.68e46}$ moleculen,
$\num{9.36e46}$ waterstofatomen.

\textbf{11.} $\num{9.36e46}/\num{9.5e27} = \qty{9.9e18}{s}$, ongeveer 310 miljard
jaar.

\textbf{12.} $\lambda_\mathrm D = 4.49\times2.014/1.008 = 8.96$; ontsnappende fractie
$9.96\,e^{-8.96} = \num{1.28e-3}$; relatieve efficiëntie $0.707\times
\num{1.28e-3}/0.0618 = 0.015$.

\textbf{13.} $H_a = kT/m_a g = \num{1.381e-23}\times200/(44.01\times\num{1.661e-27}
\times9.03) = \qty{4.18}{km}$.

\textbf{14.} $b = \qty{1.2e21}{m^{-1}.s^{-1}}$, dus
\[
\Phi_\mathrm{DL} = \frac{b f(m_a - m_i)}{m_a H_a}
= \frac{\num{1.2e21}\times\num{1e-3}\times(42.0/44.0)}{4185}
= \qty{2.74e14}{m^{-2}.s^{-1}}.
\]

\textbf{15.} $\num{3.0e13}\times\num{1e-3}\times9.03/9.81 =
\qty{2.76e10}{cm^{-2}.s^{-1}} = \qty{2.76e14}{m^{-2}.s^{-1}}$: hetzelfde
binnen 1\,\%.

\textbf{16.} $4\pi(\num{6.2e6})^2 = \qty{4.83e14}{m^2}$; $\num{1.32e29}$ moleculen
van $\ce{H2}$, $\num{2.64e29}$ atomen per seconde.

\textbf{17.} De diffusiegrens levert 28 keer zoveel als wat de \omterm{def:b3:planetary-atmospheres:exobase}{exobasis} bij de huidige
dichtheid laat ontsnappen. Waterstof hoopt zich op in de thermosfeer totdat de
Jeans-flux gelijk is aan de toevoer: op de lange termijn is de ontsnapping
diffusiebeperkt, zoals op de Aarde.

\textbf{18.} $\num{2e11}\times28 = \qty{5.6e12}{m^{-3}}$.

\textbf{19.} $\num{4.68e46}/\num{1.32e29} = \qty{3.54e17}{s}$, 11.2 miljard jaar.

\textbf{20.} $S = 0.30/0.60^2 = 0.833$.

\textbf{21.} $d = \sqrt{L/S_\mathrm{eff}}$: binnenrand $\sqrt{0.30/0.939} =
\qty{0.565}{au}$, buitenrand $\sqrt{0.30/0.290} = \qty{1.02}{au}$. Kepos,
op \qty{0.60}{au}, ligt binnen, dicht bij de binnenrand.

\textbf{22.} $0.939/0.833 = 1.13$: een toename van de helderheid met 13\,\%.

\textbf{23.} $0.833\times1361\times0.7/4 = \qty{198}{W.m^{-2}}$, $T_\mathrm{eq} =
(198/\sigma)^{1/4} = \qty{243}{K}$; ruim onder de
\qty{291}{W.m^{-2}}-grens (berekend voor een Zon-achtige ster, maar de marge
is ruim).

\textbf{24.} Drie keer zo snel: 3.7 miljard jaar.

\textbf{25.} \textbf{11~miljard~jaar}.
\end{solution}
